
%%%%%%%%%%  scan idx 370  |  folio 363  %%%%%%%%%%


S G A  7


\underline{E X P O S E  XXI}


\underline{\underline{LE NIVEAU DE LA COHOMOLOGIE DES INTERSECTIONS COMPLETES}}

\underline{par N. Katz}

avec un Appendice \quad par P. Deligne ($\S$ 5)


\underline{SOMMAIRE}

0.  Introduction. \hfill 1

1.  La matrice de Hasse-Witt d'une intersection complète. \hfill 2

2.  Le coniveau dans le cas d'un corps de définition fini. \hfill 14

3.  Application aux intersections complètes. \hfill 16

4.  Les conclusions. \hfill 19

5.  Théorèmes d'intégralité. \hfill 21


%%%%%%%%%%  scan idx 371  |  folio 364  %%%%%%%%%%


0. \underline{Introduction}

Soient $X$ une variété projective et lisse de dimension $n \geq 3$, $Y(d)$ une section par une hypersurface de degré $d$, et $E^{n-1}(Y(d),\mathbb{Q}_\ell)$ la cohomologie ``évanescente'' de $Y(d)$ (XVIII 5.4.4). Le but du présent exposé et de démontrer que pour $d$ suffisamment grand, et $Y(d)$ suffisament générale, il n'est pas vrai que tout $E^{n-1}(Y(d),\mathbb{Q}_\ell)$ ``provienne'' du $H^i(i < n-1)$ de variétés algébriques propres et lisses, par des correspondances algébriques. En caractéristique nulle, c'est une conséquence triviale de la bigraduation de Hodge sur la cohomologie complexe, et de ce que, pour la cohomologie cohérente, on a

\begin{equation*}
\dim H^{n-1}(Y(d),\ \mathcal{O}_{Y(d)}) \longrightarrow \infty \ \mbox{pour} \ d \longrightarrow \infty \quad .
\end{equation*}
La méthode employée en caractéristique $p$ est encore de relier la cohomologie cohérente $H^{n-1}(Y(d),\ \mathcal{O}_{Y(d)})$ (un vectoriel en car. $p$) à la cohomologie étale $H^{n-1}_{\mathrm{et}}(Y(d),\mathbb{Q}_\ell)$. Pour ceci, on se place sur un corps
%%%%%%%%%%  scan idx 372  |  folio 365  %%%%%%%%%%
de base fini, où les deux sont liés par la fonction $\zeta$ de la variété $Y(d)$, fonction qui s'exprime en termes de la cohomologie $\ell$-adique, et dont la ``réduction modulo $p$'' s'exprime en termes de la cohomologie cohérente (cf. XXII 3.9). Cette méthode marche sans problème pour les intersections complètes, grâce à la connaissance détaillée des deux espèces de cohomologie (étale et cohérente) dans ce cas, et donne le résultat principal 4.2 comme conséquence d'un résultat (1.4) sur les opérations de Hasse-Witt (1.0).

Il existe des arguments heuristiques dûs à A. GROTHENDIECK et s'appuyant sur le yoga de la cohomologie cristalline, qui rendent plausible l'énoncé général pour toute $X$ projective et lisse, par essentiellement la même méthode.

Même dans le cas particulier des intersections complètes, la démonstration utilise de façon essentielle (2.2) le théorème d'intégralité des valeurs propres de Frobenius 2.1, démontré récemment par P. DELIGNE, et qui est prouvé dans l'Appendice I.


1. \underline{La matrice de Hasse-Witt d'une intersection complète}

\underline{Définition} 1.0. \underline{Un schéma propre sur un corps} $k$ \underline{d'exposant caractéristique} $p$ \underline{s'appelle} spécial \underline{si l'opération additive} $p$-\underline{linéaire}

\begin{equation}\tag{1.0.1}
\mathfrak{F}_p : \mathcal{O}_X \longrightarrow \mathcal{O}_X \quad , \quad g \longmapsto g^p
\end{equation}
\underline{induit une opération}

\begin{equation}\tag{1.0.2}
\mathfrak{F}_p : H^{\dim X}(X,\mathcal{O}_X) \longrightarrow H^{\dim X}(X,\mathcal{O}_X) 
\end{equation}

%%%%%%%%%%  scan idx 373  |  folio 366  %%%%%%%%%%

 \underline{qui est nilpotente (ou, ce qui revient au même, posant} $M = \dim_k(H^{\dim X}(X,\mathcal{O}_X))$, \underline{si l'opération $p^M$-linéaire}

\begin{equation}\tag{1.0.3}
\mathfrak{F}_p^M = \mathfrak{F}_{p^M} : \mathcal{O}_X \longrightarrow \mathcal{O}_X \quad , \quad g \longmapsto g^{p^M}
\end{equation}
\underline{tue} $H^{\dim X}(X,\mathcal{O}_X)$.

Rappelons [2] que pour tout $i \in \mathbb{Z}$ , l'homomorphisme

\begin{equation*}
H^i(X,\mathcal{O}_X) \stackrel{\mathfrak{F}_p}{\longrightarrow} H^i(X,\mathcal{O}_X)
\end{equation*}
induit par (1.0.1) s'appelle classiquement \underline{l'homomorphisme de Hasse-Witt} dans $H^i(X,\mathcal{O}_X)$. On notera que si $p = 1$, dire que $X$ est non-spécial signifie simplement que $H^{\dim X}(X,\mathcal{O}_X) \neq 0$.

\underline{Théorème 1.1.} \underline{Fixons} \underline{un exposant caractéristique} $p$, \underline{une dimension} $n \geq 1$, \underline{un entier} $r \geq 1$ \underline{et un multidegré}$(d_1,\dots,d_r)$. \underline{Si on a}

\begin{equation}\tag{1.1.0}
\sum_{i=1}^{r} d_i \geq (n+r)+1
\end{equation}
\underline{alors, parmi toute les intersections complètes} (\underline{de dimension} $n$) \underline{dans} $\mathbb{P}^{n+r}$ \underline{de multidegré} $(d_1,\dots,d_r)$ \underline{en caractéristique} $p$, \underline{un élément} \underline{``assez général'' est non-spécial} (1.0).

\underline{Démonstration.} Désignons par

\begin{equation}\tag{1.1.1}
f : X_U \longrightarrow U
\end{equation}
la``famille universelle'' des intersections complètes dans $\mathbb{P}^{n+r}$ de
%%%%%%%%%%  scan idx 374  |  folio 367  %%%%%%%%%%
multidegré $(d_1,\dots,d_r)$ en caractéristique $p$ (muni des équations de définition), de sorte que $U$ est un ouvert dans un produit de $r$ espaces projectifs sur $\mathbb{F}_p$ , d'une dimension impressionnante.

On sait que le $\mathcal{O}_U$-module

\begin{equation}\tag{1.1.2}
R^n f_* \mathcal{O}_{X_U}
\end{equation}
est localement libre, et sa formation commute à tout changement de base (XI 1.5 (i)) et, par l'hypothèse (1.1.0), que son rang $M$ est non nul (XI 2.5). Cela donne le théroème si $p = 1$, pour toute intersection complète vérifiante (1.1.0).

Supposons $p \geq 2$. L'opération de Hasse-Witt itérée

\begin{equation}\tag{1.1.3}
\mathfrak{F}_p^M = \mathfrak{F}_{p^M} : R^n f_* \mathcal{O}_{X_U} \longrightarrow R^n f_* \mathcal{O}_{X_U}
\end{equation}
étant $p^M$-linéaire sur $\mathcal{O}_U$ , et $R^n f_* \mathcal{O}_{X_U}$ étant localement libre, l'action (1.1.3) s'exprime localement sur $U$ par une matrice $H$ de type $(M,M)$ à valeurs dans $\mathcal{O}_U$ , de sorte que

\begin{equation}\tag{1.1.4}
\{ u \in U \mid X_u \mbox{ est non-spéciale}\} \mbox{ est ouvert dans } U
\end{equation}
(car c'est l'ensemble des $u$ tels que $H(u) \neq 0$), et, $U$ étant irréductible, il suffit de trouver une valeur de $u$ telle que $X_u$ soit non spéciale, ce qu'on va faire sur un corps fini.

\underline{Proposition 1.1.5.} \underline{Soit} $X$ \underline{une intersection complète sur} $\mathbb{F}_q$. \underline{Si} $X$ \underline{n'a pas des points} $\mathbb{F}_q$\underline{-rationnels,} \underline{alors} $X$ \underline{est non spéciale.}


%%%%%%%%%%  scan idx 375  |  folio 368  %%%%%%%%%%

\underline{Démonstration.} En effet, il suffit évidemment de prouver que pour $X$ une intersection complète, on a l'implication

\begin{equation}\tag{1.1.5.1}
X \text{ spéciale} \Longrightarrow \mathrm{card}(X(\mathbb{F}_q)) \equiv 1 \bmod p \quad .
\end{equation}
Or, d'après la formule de la trace (prouvée dans l'exposé suivant XXII, 3.2), on a

\begin{equation*}
\mathrm{card}(X(\mathbb{F}_q)) \equiv \sum_{i \geq 0} (-1)^i \, \mathrm{tr}(\mathfrak{F}_q \mid H^i(X,\mathcal{O}_X)) \quad .
\end{equation*}
Compte tenu de ce que $H^i(X,\mathcal{O}_X)=0$ si $i \neq 0, n=\dim(X)$ et $H^0(X,\mathcal{O}_X) \simeq \mathbb{F}_q$ (XI, 1.5 (iii)), ceci se récrit

\begin{equation*}
\mathrm{card}(X(\mathbb{F}_q)) \equiv 1 + (-1)^n \, \mathrm{tr}(\mathfrak{F}_q \mid H^n(X,\mathcal{O}_X)) \quad .
\end{equation*}
Si $X$ est spéciale, alors par définition l'opération de $\mathfrak{F}_q$ sur $H^n(X,\mathcal{O}_X)$ est nilpotente, donc a une trace nulle, d'où l'implication (1.1.5.1) et la proposition.

\underline{Proposition} 1.2. \underline{Quelle que soit la puissance} $q$ \underline{de} $p$, \underline{si}

\begin{equation*}
\sum_1^r d_i \geq n+r+1 \quad ,
\end{equation*}
\underline{il existe une intersection complète dans} $\mathbb{P}^{n+r}_{\mathbb{F}_q}$ \underline{de multidegré} $(d_1,\dots,d_r)$ \underline{qui est sans point} $\mathbb{F}_q$\underline{-rationnel}.

\underline{Démonstration.}

Commençons par le cas $r=1$. Pour tout $d \geq n+2$, il faut trouver une forme homogène de degré $d$ 


%%%%%%%%%%  scan idx 376  |  folio 369  %%%%%%%%%%

\begin{equation}\tag{1.2.1}
 H(X_1),\dots,X_{n+2}) \in \mathbb{F}_q[X_1,\dots,X_{n+2}]
\end{equation}
telle que

\begin{equation}\tag{1.2.2}
(\alpha_1,\dots,\alpha_{n+2}) \in (\mathbb{F}_q)^{n+2}, \; H\,(\alpha_1,\dots,\alpha_{n+2}) = 0 \Longrightarrow \alpha_1 = \dots = \alpha_{n+2} = 0 \quad .
\end{equation}
Pour ceci, posons $\mathbb{E} = \mathbb{F}_{q^d}$ , et soient $\beta_1,\dots,\beta_{n+2}$ des éléments de $\mathbb{E}$ qui sont linéairement indépendant sur $\mathbb{F}_q$ . On prend

\begin{equation}\tag{1.2.3}
H(X_1,\dots,X_{n+2}) \;=\; N_{\mathbb{E}/\mathbb{F}_q} \Bigl( \sum_{i=1}^{n+2} \beta_i X_i \Bigr) \quad .
\end{equation}
Dans le cas général, compte tenu de ce que

\begin{equation}\tag{1.2.4}
\sum_i^r d_i \geq n+r+1 \quad ,
\end{equation}
on peut trouver $r$ entiers

\begin{equation}\tag{1.2.5}
1 \leq a_1 < a_2 < \dots < a_r = n+r+1
\end{equation}
vérifiant

\begin{equation}\tag{1.2.6}
\left\{ \begin{array}{lcl} d_1 & \geq & a_o \\[3ex] d_2 & \geq & a_2 - a_1 \\[3ex] d_r & \geq & a_r - a_{r-1} \end{array} \right. \quad .
\end{equation}

%%%%%%%%%%  scan idx 377  |  folio 370  %%%%%%%%%%

On prend, pour $i = 1,\dots,r$, une forme homogène de degré $d_i$

\begin{equation*}
H_i \, (X_{(a_{i-1}+1)},\dots,X_{a_i}) \ \in \ \mathbb{F}_q[X_{a_{i-1}+1},\dots,X_{a_i}]
\end{equation*}
en $a_i - a_{i-1}$ variable  qui a la propriété (1.2.2). L'intersection complète voulue est celle définie par les formes homogènes

\begin{equation*}
H_1(X_1,\dots,X_{a_1}), \ \ \dots\ , \ \ H_r(X_{a_{r-1}},\dots,X_{a_r}) \qquad .
\end{equation*}
\underline{Théorème} 1.3.  \underline{Soit} $X$ \underline{une intersection complète}, \underline{irréductible}, \underline{de dimension} $n \geq 1$ \underline{sur un corps} $k$ \underline{d'exposant caractéristique} $p$. \underline{Pour tout entier} $d \geq 1 + n$, \underline{une section assez générale de} $X$ \underline{par une hypersurface de degré} $d$ \underline{est non-spéciale} (1.0).

\underline{Démonstration}. Les sections $X.H_u$ de $X$ (par des hypersurfaces) de degré $d$ sont (paramétérisées par) les droites dans $H^0(X,\mathcal{O}_X(d))$, i.e. par

\begin{equation}\tag{1.3.1}
\mathbb{P}(H^0(X,\mathcal{O}_X(d)) = U
\end{equation}
et au-dessus de $U$, nous avons la ``section universelle de $X$ par une hypersurface de degré $d$ ''

\begin{equation}\tag{1.3.2}
f : X.H_u \longrightarrow U \qquad ,
\end{equation}
soit $U'$ l'ouvert des $u \in U$ tels que $\dim(X_u) = n-1$. Nous savons par (XI 1.5 (i) et 2.5) que le $\mathcal{O}_U$-module

\begin{equation}\tag{1.3.3}
R^{n-1}f_*(\mathcal{O}_{X.H_u}) \ \mid \ U' 
\end{equation}

%%%%%%%%%%  scan idx 378  |  folio 371  %%%%%%%%%%

 est un faisceau localement libre dont la formation commute à tout changement de base, et est de rang $M \geq 1$. Cela donne le théorème si $p = 1$, pour toute section par une hypersurface de degré $\geq 1+n$ .

On suppose donc que $p > 0$ ; l'opération $\mathfrak{J}_{p^n}$ s'exprime encore localement en termes d'une matrice $M \times M$ à valeurs dans $\mathcal{O}_u$, de sorte que

\begin{equation}\tag{1.3.4}
\{ \, u \in U' | X.H_u \text{ est non-spéciale} \} \ \text{ est ouvert dans } U' \ .
\end{equation}
Il suffit donc de démontrer que l'ensemble (1.3.4) est \underline{non-vide}.

Traitons d'abord le cas

\begin{equation}\tag{1.3.5}
k = \mathbb{F}_q \qquad\qquad .
\end{equation}
Quitte à remplacer $\mathbb{F}_q$ par une extension finie, on peut supposer

\begin{equation}\tag{1.3.6}
\mathrm{card}(X(\mathbb{F}_q)) \ > \ 0
\end{equation}
\begin{equation}\tag{1.3.7}
\begin{aligned}
&\text{il existe des coordonnées homogènes } X_1,\dots,X_{n+r+1} \\
&\text{dans l'espace ambiant } \mathbb{P}^{n+r}_k, \text{ telles que } X \text{ ne rencontre pas} \\
&\text{la variété linéaire } X_1=\dots=X_{n+1}=0.
\end{aligned}
\end{equation}
On prend alors une forme homogène de degré $d$

\begin{equation}\tag{1.3.8}
H(X_1,\dots,X_{n+1}) \ \in \ \mathbb{F}_q[X_1,\dots, X_{n+1}]
\end{equation}
qui a la propriété (1.2.2) (ceci est possible car $d \geq 1+n$). On prend pour $H$ l'hypersurface dans $\mathbb{P}^{n+r}$ d'équation $H(X_1,\dots,X_{n+1}) = 0$. Alors
%%%%%%%%%%  scan idx 379  |  folio 372  %%%%%%%%%%
l'intersection $X.H$ est sans point $\mathbb{F}_q$-rationnel par construction, ce qui montre à la fois, grâce à (1.3.6), que $X \not\subset H$, donc $X.H$ est de dimension $n-1$ ($X$ étant irréductible) i.e. $H \in U'$, et que $X.H$ est non-spécial (1.5).

Passons au cas général. Regardant les équations qui définissent $X$, on voit qu'on peut supposer que $k$ est de type fini sur $\mathbb{F}_p$, et qu'il existe un sous-anneau $B$ de $k$, de type fini (comme anneau) sur $\mathbb{F}_p$ dont le corps des fractions est $k$, tel que (1.3.9) $X \subset \mathbb{P}^{n+r}_k$ est la fibre générique d'une intersection complète relative $\widetilde{X} \subset \mathbb{P}^{n+r}_B$ qui est, \underline{fibre par fibre}, de même multidegré que $X$, et géométriquement irréductible.

Désignons par

\begin{equation}\tag{1.3.10}
\rho : \widetilde{X} \longrightarrow S = \mathrm{Spec}(B)
\end{equation}
la projection, par

\begin{equation}\tag{1.3.11}
E, \ \mbox{le } \mathcal{O}_S\mbox{-module localement libre } \rho_*\mathcal{O}_{\widetilde{X}}(d) \ ,
\end{equation}
et par

\begin{equation}\tag{1.3.12}
\pi : T = \check{\mathbb{P}}(E) \longrightarrow S
\end{equation}
le fibré associé des droites dans $\mathbb{E}$ . Notons que la fibre générique de (1.3.12) n'est autre que $U = \mathbb{P}(H^o(X,\mathcal{O}_X(d))$. Au-dessus de $T$ nous avons la section universelle de $\widetilde{X} \subset \mathbb{P}^{n+r}_B$ par des hypersurfaces de degré $d$, notée 


%%%%%%%%%%  scan idx 380  |  folio 373  %%%%%%%%%%

\begin{equation}\tag{1.3.13}
 \alpha : \widetilde{X}.H_T \longrightarrow T \ .
\end{equation}
Ces constructions s'insèrent dans le diagramme à carrés cartésiens

\begin{equation}\tag{1.3.14}
\begin{tikzcd}
\widetilde{X}.H_T \arrow[d, "\alpha"] & X.H_u = (\widetilde{X}.H_T) \times_U T \arrow[l] \arrow[d] \\
T \arrow[d, "\pi"] & U = T \times_S \mathrm{Sp}(k) \arrow[l] \arrow[d] \\
S & \mathrm{Sp}(k) \arrow[l]
\end{tikzcd}
\end{equation}
% STRUCTURE: 6 nodes, arranged in a 3-row by 2-column grid. NODES: row 1 left = $\widetilde{X}.H_T$ ;
%   row 1 right = $X.H_u = (\widetilde{X}.H_T) \times_U T$ ; row 2 left = $T$ ; row 2 right = $U = T
%   \times_S \mathrm{Sp}(k)$ ; row 3 left = $S$ ; row 3 right = $\mathrm{Sp}(k)$. ARROWS: 7 in total,
%   every one a plain solid single-shaft arrow with a single open V-head — none is hooked/mono, none is
%   a two-headed epi, none carries an iso or equality label, none is dashed or dotted, and no
%   pullback/cartesian corner marks are drawn. (1) horizontal, from row 1 right node $X.H_u =
%   (\widetilde{X}.H_T) \times_U T$ to row 1 left node $\widetilde{X}.H_T$, pointing LEFT, unlabelled.
%   (2) vertical, from row 1 left node $\widetilde{X}.H_T$ to row 2 left node $T$, pointing DOWN,
%   labelled $\alpha$, the label sitting to the RIGHT of the shaft, about mid-height. (3) vertical, from
%   row 1 right node $X.H_u = (\widetilde{X}.H_T) \times_U T$ to row 2 right node $U = T \times_S
%   \mathrm{Sp}(k)$, pointing DOWN, unlabelled. (4) horizontal, from row 2 right node $U = T \times_S
%   \mathrm{Sp}(k)$ to row 2 left node $T$, pointing LEFT, unlabelled. (5) vertical, from row 2 left
%   node $T$ to row 3 left node $S$, pointing DOWN, labelled $\pi$, the label sitting to the RIGHT of
%   the shaft, about mid-height. (6) vertical, from row 2 right node $U = T \times_S \mathrm{Sp}(k)$ to
%   row 3 right node $\mathrm{Sp}(k)$, pointing DOWN, unlabelled. (7) horizontal, from row 3 right node
%   $\mathrm{Sp}(k)$ to row 3 left node $S$, pointing LEFT, unlabelled. The tag (1.3.14) is printed in
%   the left margin level with row 2. A comma is printed on the far right of the row-3 line, well clear
%   of the diagram; it is not a node and is carried at the head of the following prose block.
, et nous avons encore que le $\mathcal{O}_T$-module

\begin{equation}\tag{1.3.15}
R^{n-1}\alpha_*(\mathcal{O}_{\widetilde{X}.H_T})
\end{equation}
au-dessus de l'ouvert $T' \subset T$ formé des $t$ tels que $\dim \alpha^{-1}(t) = n-1$, est localement libre, et se formation y commute à tout changement de base. Compte tenu de ce que les schémas $T$ et $U$ ont le \underline{même} point générique, pour démontrer que (1.3.4) est non vide, il suffit de vérifier que l'ouvert de $T'$

\begin{equation}\tag{1.3.16}
\{\ t \in T' \mid \widetilde{X}_{\pi(t)}.H_t \ \mbox{est non-spéciale}\}
\end{equation}
est non vide. Or, le cas (1.3.5) déjà traité démontre que pour \underline{tout point}
%%%%%%%%%%  scan idx 381  |  folio 374  %%%%%%%%%%
\underline{fermé} $s \in S$, la fibre $\pi^{-1}(s)$ a une intersection non vide avec l'ouvert (1.3.16), Q.E.D.

Des raisonnements tout à fait semblables permettent de démontrer:

\underline{Théorème} 1.4. \underline{Soit} $X \subset \mathbb{P}^N$ \underline{un schéma projectif de Cohen-Macaulay,} \underline{équidimensionnel de dimension} $n \geq 1$, \underline{sur un corps} $k$, \underline{alors il existe} \underline{un entier} $d_o$ \underline{tel que pour} $d \geq d_o$, \underline{désignant par} $X.H \stackrel{i}{\longrightarrow} X$ \underline{l'inclusion} \underline{d'une section par une hypersurface} générale \underline{de degré} $d$, \underline{l'action de}

\begin{equation}\tag{1.4.0}
\mathfrak{F}_p \ \mathrm{sur}\ H^{n-1}(X.H,\mathcal{O}_{X.H}) \ /\ i^{*}H^{n-1}(X,\mathcal{O}_X)
\end{equation}
\underline{soit} non nilpotente.

\underline{Démonstration.} On se ramène, tout comme ci-dessus, au cas $k = \mathbb{F}_q$, et au problème de trouver un \underline{seul} $H$ (pour un degré donné assez grand et tel que $\dim(H^{n-1}(X.H,\mathcal{O}_{X.H}))$ soit indépendant de $H$ ; cf. 1.4.4) qui rend vrai (1.4.0). Quitte à faire agrandir $q$ , on peut supposer que les composantes irréductibles $X_i$ de $X$ sont géométriquement irréductibles, satisfont

\begin{equation}\tag{1.4.1}
\mathrm{Card}\ X_i(\ \mathbb{F}_q) \geq 2 \quad ,
\end{equation}
et qu'on ait

(1.4.2) il y a des coordonnés projectives $X_o,\ldots,X_N$ telles que $X$ ne coupe pas la variété linéaire $X_o = \ldots = X_n = 0$ .

Par un théorème d'évanescence de Serre, ([1] ou SGA 2 XII 1.4), il existe un entier $d_1$ tel que pour tout entier $d \geq d_1$ , on ait :


%%%%%%%%%%  scan idx 382  |  folio 375  %%%%%%%%%%

\begin{equation}\tag{1.4.3}
H^i(X,\mathcal{O}_X(-d)) = 0 \ \mathrm{si}\ i \neq n = \dim X \quad .
\end{equation}
Désignant par $H$ une hypersurface de degré $d \geq d_1$, on a (par 1.4.3) :

\begin{equation}\tag{1.4.4}
\begin{aligned} H^i(X,\mathcal{O}_X) &\stackrel{\sim}{\longrightarrow} H^i(X.H,\mathcal{O}_{X.H}) \qquad \mathrm{si}\ \ i \leq n - 2 \\[1ex] H^{n-1}(X,\mathcal{O}_X) &\lhook\joinrel\longrightarrow H^{n-1}(X.H,\mathcal{O}_{X.H})\ \mathrm{injectif}\ . \end{aligned}
\end{equation}
Nous allons démontrer le théorème en prenant $d_o = \max\ (1+N,d_1)$. Pour tout entier $d \geq d_o$, il faut trouver une hypersurface $H$ de degré $d$, définie sur $\mathbb{F}_q$, qui ne contienne aucun $X_i$, et telle que

\begin{equation}\tag{1.4.5}
\mathrm{trace}(\mathfrak{F}_q\ \mathrm{sur}\ H^{n-1}(X.H,\mathcal{O}_{X.H})\ /\ H^{n-1}(X,\mathcal{O}_X)) \neq 0 \quad .
\end{equation}
D'après la formule de la trace (XX, 3.2.1), on a

\begin{equation}\tag{1.4.6}
\begin{aligned}
\mathrm{card}(X.H\ (\mathbb{F}_q))\mathrm{mod}\ p
&\equiv \sum_{i=0}^{n-1}(-1)^i
\mathrm{tr}(\mathfrak{F}_q\ \mathrm{sur}\ H^i(X.H,\mathcal{O}_{X.H})) \quad , \\
&\stackrel{(1.4.4)}{\equiv} \sum_{i=0}^{n-1}(-1)^i
\mathrm{tr}(\mathfrak{F}_q\ \mathrm{sur}\ H^i(X,\mathcal{O}_X)) \\
&\qquad +(-1)^{n-1}\mathrm{tr}\!\left(\mathfrak{F}_q\ \mathrm{sur}\
\frac{H^{n-1}(X.H,\mathcal{O}_{X.H})}{H^{n-1}(X,\mathcal{O}_X)}\right) .
\end{aligned}
\end{equation}
Désignons par $a \in \mathbb{Z}/p\ \mathbb{Z}$ , la quantité

\begin{equation*}
\sum_{i=0}^{n-1} (-1)^i \mathrm{tr}(\mathfrak{F}_q\ \mathrm{sur}\ H^i(X,\mathcal{O}_X)) \quad .
\end{equation*}

%%%%%%%%%%  scan idx 383  |  folio 376  %%%%%%%%%%

Il suffit alors, pour tout $d \geq d_o$, de trouver une $H$ défini sur $\mathbb{F}_q$, de degré $d$, qui ne contient pas $X$, et tel que

\begin{equation}\tag{1.4.7}
\mathrm{Card}(X.H\ (\mathbb{F}_q)) \not\equiv a \bmod p \ .
\end{equation}
Il y a lieu de distinguer deux cas. Si $a \neq 0$, il suffit de prendre un $H$ sur $\mathbb{F}_q$ tel que $(X.H)\ (\mathbb{F}_q)$ soit vide, tout comme dans (1.3.8). Un tel $H$ vérifie (1.4.7), et ne contient aucun $X_i$ (par hypothèse on sait que $\mathrm{card}(X_i(\mathbb{F}_q)) \geq 2$).

Si $a = 0$, il suffira de trouver un $H$ tel que $\mathrm{Card}(X.H)(\mathbb{F}_q)=1$ (1.4.7 sera vérifié, et $X_i \not\subset H$, car $\mathrm{Card}(X_i)(\mathbb{F}_q) \geq 2$). Pour ceci, on choisit des coordonnées projectives $X_o,\dots,X_N$ telles que $(1,0,\dots,0)$ soit un point de $X$. On prend $H$ d'équation $H(X_1,\dots,X_N) = 0$, où $H$ est une forme homogène telle que (cf. (1.2.2))

\begin{equation}\tag{1.4.8}
\alpha_1,\dots,\alpha_N \in \mathbb{F}_q, \ H(\alpha_1,..,\alpha_N) = 0 \Longrightarrow \alpha_1=\dots=\alpha_N = 0 \ .
\end{equation}
Alors, l'hypersurface définie par $H$ dans $\mathbb{P}^N$ n'a que le point $(1,0,\dots,0)$ comme point à valeurs dans $\mathbb{F}_q$, donc a fortiori, $X.H$ n'a que ce point comme $\mathbb{F}_q$-point. cqfd.


%%%%%%%%%%  scan idx 384  |  folio 377  %%%%%%%%%%


2. \underline{Le coniveau dans le cas d'un corps de définition fini}

2.0. Soit $X$ un schéma séparé de type fini sur $\mathbb{F}_q$. Pour tout nombre premier $\ell \neq \mathrm{car}(\mathbb{F}_q)$, on désigne par

\begin{equation}\tag{2.0.1}
H^i_c(\overline{X},\mathbb{Q}_\ell)
\end{equation}
les groupes de cohomologie à supports propres de $\overline{X} = X \otimes_{\mathbb{F}_q} \overline{\mathbb{F}}_q$, qui sont des $\mathbb{Q}_\ell$-vectoriels de rang fini sur lesquels opère le groupe de Galois $\mathrm{Gal}(\overline{\mathbb{F}}_q/\mathbb{F}_q)$ de $\mathbb{F}_q$ . Désignons par

\begin{equation}\tag{2.0.2}
\varphi_q \in \mathrm{Gal}(\overline{\mathbb{F}}_q/\mathbb{F}_q) \quad \mbox{l'automorphisme de Frobenius} \ .
\end{equation}
On sait (SGA 5 XIII 3.2.1) que la fonction zêta de $X$ (XXII 3.0.1) est donné par la formule

\begin{equation}\tag{2.0.3}
Z(t,X/\mathbb{F}_q) = \prod_{i \geq 0} \det (1-t\ \varphi_q^{-1}|H^i_c(\overline{X},\mathbb{Q}_\ell))^{(-1)^{i+1}} \ .
\end{equation}
DELIGNE a démontré récemment, en généralisant un résultat dû à WASHNITZER dans le cas $X$ projectif et lisse, le théorème suivant, qui est prouvé dans l'appendice plus bas (5.5.3) :

\underline{Théorème} 2.1. \underline{Les valeurs propres de} $\varphi_q^{-1}$ \underline{dans} $H^i_c(\overline{X},\mathbb{Q}_\ell)$ \underline{sont des entiers algébriques}, \underline{et}, \underline{si} $i > \dim X$, \underline{elles sont divisibles} (\underline{en tant qu'entiers algébriques}) \underline{par} $q^{i-\dim X}$ .


%%%%%%%%%%  scan idx 385  |  folio 378  %%%%%%%%%%

En particulier, on obtient (tautologiquement),

(2.1 bis) \quad les valeurs propres de $\varphi_q^{-1}$ dans $H^i_c(\overline{X},\mathbb{Q}_\ell(-j))$ sont des entiers algébriques divisibles par $q^j$ .

\underline{Corollaire} 2.2. \underline{Soit} $X$ \underline{propre lisse et géométriquement connexe sur} $\mathbb{F}_q$ . \underline{Alors les valeurs propres de} $\varphi_q^{-1}$ \underline{dans} $N^jH^q(\overline{X},\mathbb{Q}_\ell)$ \underline{(notations de} XX 2.0.6) \underline{sont des entiers algébriques divisibles par} $q^j$ .

\underline{Démonstration.} L'énoncé est invariant par extension finie $\mathbb{F}_{q^M}/\mathbb{F}_q$ (qui revient à prendre les puissances. $M$\textsuperscript{èmes} des valeurs propres envisagées). Quitte à faire un tel changement de base, on a, par définition un nombre fini de schémas $T_i$ projectifs, lisses, et géométriquement connexes sur $\mathbb{F}_q$, des dimensions $d_i$, et de cycles algébriques

$Z_i$ sur $T_i \times_{\mathbb{F}_q} X$ de codimension $j + d_i$

tels que \underline{l'homomorphisme de Gal}$(/\mathbb{F}_q)$-\underline{modules}

\begin{equation*}
H^{q-2j}(\overline{T}_i,\ \mathbb{Q}_\ell(-j)) \ \stackrel{\oplus \mathrm{Cl}(Z_i)}{\longrightarrow} \ N^jH^q(\overline{X},\mathbb{Q}_\ell)
\end{equation*}
soit surjectif, et on conclut par 2.1 bis.


%%%%%%%%%%  scan idx 386  |  folio 379  %%%%%%%%%%


3. \underline{Application aux intersections complètes}

3.0. Soit $X$ une intersection complète dans $\mathbb{P}^N$ lisse de dimension $n \geq 1$ sur un corps $k$. Pour $\ell \neq \mathrm{car}(k)$, on sait par ``Lefschetz faible'' et la dualité que (XI 1.6)

\begin{equation}\tag{3.0.1}
H^i(\overline{X},\mathbb{Q}_\ell) \;=\; \left\{ \begin{array}{l} 0 \ \mbox{si $i$ impair et $i \neq n$} \\[1.5ex] \mathbb{Q}_\ell\!\left(\frac{-i}{2}\right) \ \mbox{si $i$ pair $\leq 2n$ et $i \neq n$} \\[1.5ex] H^n(\overline{X},\mathbb{Q}_\ell) \ \mbox{si $i=n$} \quad . \end{array} \right.
\end{equation}
Pour $k = \mathbb{F}_q$, on obtient donc une formule pour la fonction zêta :

\begin{equation}\tag{3.0.2}
Z(t,X/\mathbb{F}_q) \ \prod_{i=0}^{n} \ (1-q^i t) = \det(1-t\,\varphi_q^{-1} \mid \mathrm{Prim}^n(\overline{X},\mathbb{Q}_\ell))^{(-1)^{n+1}} \quad ,
\end{equation}
(on a employé da décomposition primitive, et (3.0.1).)

\underline{Corollaire} 3.0.3. \underline{Sous les hypothèses} 3.0, \underline{avec} $k = \mathbb{F}_q$ , \underline{on a une congruence modulo $p$}

\begin{equation*}
\det(1-t\,\varphi_q^{-1} \mid H^n(\overline{X},\mathbb{Q}_\ell)) \;\equiv\; \det(1-t\,\mathfrak{F}_q \mid H^n(X,\mathcal{O}_X)) \quad .
\end{equation*}
\underline{Démonstration.} D'après la formule de congruence (XXII, 3.2.1), et compte tenu de ce que

\begin{equation}\tag{3.0.4}
H^i(X,\mathcal{O}_X) \;\simeq\; \left\{ \begin{array}{ll} \mathbb{F}_q & i = 0 \\[1.5ex] 0 & 1 \leq \; \leq n - 1 \\[1.5ex] 0 & i \geq n+1 \end{array} \right.
\end{equation}

%%%%%%%%%%  scan idx 387  |  folio 380  %%%%%%%%%%
on a

\begin{equation}\tag{3.0.5}
Z(t,X/\mathbb{F}_q) \;\equiv\; (1-t)^{-1}\ \det\,(1-t\mathfrak{F}_q|H^n(X,\mathcal{O}_X))^{(-1)^{n+1}} \quad \mathrm{mod}\ p
\end{equation}
de sorte que (3.0.5) et (3.0.2) donnent la congruence

\begin{equation}\tag{3.0.6}
\det(1-t\ \varphi_q^{-1}|\mathrm{Prim}^n(\overline{X},\mathbb{Q}_\ell)) \;\equiv\; \det(1-t\ \mathfrak{F}_q|H^n(X,\mathcal{O}_X)) \quad \mathrm{mod}\ p \quad .
\end{equation}
La conclusion voulue s'en déduit, compte tenu de ce que

\begin{equation}\tag{3.0.7}
H^n(\overline{X},\mathbb{Q}_\ell) \;\simeq\; \begin{cases} \mathrm{Prim}^n(\overline{X},\mathbb{Q}_\ell) & \text{si } n \text{ est impair} \\[2mm] \mathrm{Prim}^n(\overline{X},\mathbb{Q}_\ell) \oplus \mathbb{Q}_\ell\!\left(\frac{-n}{2}\right) & \text{si } n \text{ est pair} \end{cases} \quad .
\end{equation}
\underline{Corollaire} 3.1. \underline{Soit} $X$ \underline{une intersection complète}, \underline{lisse de dimension} $n \geq 1$ \underline{sur} $\mathbb{F}_q$ . \underline{Si}

\begin{equation}\tag{3.1.0}
H^n(\overline{X},\mathbb{Q}_\ell) \;=\; N^1H^n(\overline{X},\mathbb{Q}_\ell)
\end{equation}
\underline{alors}

\begin{equation}\tag{3.1.1}
\mathfrak{F}_p\ \underline{\text{est nilpotent sur}}\ H^n(X,\mathcal{O}_X) \quad .
\end{equation}
\underline{Démonstration}. D'après (2.2), (3.1.0) implique qu'on a

\begin{equation}\tag{3.1.2}
\det(1-t\ \varphi_q^{-1}\ |\ H^n(\overline{X},\mathbb{Q}_\ell)) \;\equiv\; 1 \quad \mathrm{mod}\ p\ (\text{et même mod } q) \quad ,
\end{equation}
et ceci, avec (3.0.3), implique qu'on a

\begin{equation}\tag{3.1.3}
\det(1-t\ \mathfrak{F}_q\ |\ H^n(X,\mathcal{O}_X)) \;=\; 1 \quad , 
\end{equation}

%%%%%%%%%%  scan idx 388  |  folio 381  %%%%%%%%%%

 ce qui implique que $\mathfrak{F}_q$ est nilpotent sur $H^n(X,\mathcal{O}_X)$, et donc que $\mathfrak{F}_p$ l'est, cqfd.

\underline{Proposition} 3.2. Soit $X$ une intersection complète, lisse de dimension $n \geq 1$ sur un corps $k$. Si

\begin{equation}\tag{3.2.1}
H^n(\overline{X},\mathbb{Q}_\ell) \;=\; N^1H^n(\overline{X},\mathbb{Q}_\ell)
\end{equation}
alors $X$ est spéciale, i.e.

\begin{equation}\tag{3.2.2}
\mathfrak{F}_p\ \text{est nilpotente sur}\ H^n(X,\mathcal{O}_X) \quad .
\end{equation}
\underline{Démonstration}. On peut supposer $k$ de type fini sur $\mathbb{F}_p$ (le cas $p = 0$ étant trivial), et qu'il existe une sous-$\mathbb{F}_p$-algèbre $B$ de type fini de $k$, telle que $X/k$ soit la fibre générique d'une intersection complète lisse relative $\widetilde{X}$ :

\begin{equation}\tag{3.2.3}
\begin{tikzcd}
\widetilde{X} \arrow[d, "\pi"'] & X = \widetilde{X}/_{\eta} \arrow[l] \arrow[d] \\
\mathrm{Spec}(B) = S & \eta = \mathrm{Sp}(k) \;\; . \arrow[l]
\end{tikzcd}
\end{equation}
% STRUCTURE: Commutative square, 4 nodes, 4 arrows. NODES: N1 top-left = $\widetilde{X}$ (capital X with
%   tilde accent); N2 top-right = $X = \widetilde{X}/_{\eta}$ (the eta is set low/right after the
%   solidus, as a subscript); N3 bottom-left = $\mathrm{Spec}(B) = S$; N4 bottom-right = $\eta =
%   \mathrm{Sp}(k)$, followed on the same baseline by a free-standing full stop '.' set to its right
%   (kept inside the node in the tikzcd). ARROWS: A1 horizontal top, from N2 to N1, pointing LEFT
%   (arrowhead at N1), long plain single-shafted arrow, NO label. A2 vertical left, from N1 down to N3,
%   pointing DOWN (arrowhead at N3), long plain arrow, labelled $\pi$ placed to the LEFT of the shaft at
%   mid-height. A3 vertical right, from N2 down to N4, pointing DOWN (arrowhead at N4), long plain
%   arrow, NO label. A4 horizontal bottom, from N4 to N3, pointing LEFT (arrowhead at N3), long plain
%   arrow, NO label. No hooked/mono arrows, no iso (\simeq) marks, no double-headed arrows, no dashed
%   arrows, no cartesian-square corner mark. The tag (3.2.3) is printed in the left margin at the
%   vertical mid-height of the square, level with the label $\pi$.
On sait que le $B$-module $R^n\pi_*(\mathcal{O}_{\widetilde{X}})$ est localement libre, et sa formation commute à tout changement de base, donc l'ensemble des points 


%%%%%%%%%%  scan idx 389  |  folio 382  %%%%%%%%%%

\begin{equation}\tag{3.2.4}
 \{\, s \in S \mid X_s \text{ est spéciale} \,\}
\end{equation}
est un fermé de $S$.

Par (XX, 2.3), quitte à rapetisser $S$, nous pouvons supposer que

\begin{equation}\tag{3.2.5}
\forall\, s \in S, \quad H^n(X_{\bar{s}}, \mathbb{Q}_\ell) \;=\; N^1 H^n(X_{\bar{s}}, \mathbb{Q}_\ell) \quad .
\end{equation}
Par (3.1), pour tout point \underline{fermé} $s$ de $S$, $X_s$ est spécial. Il s'ensuit que le fermé (3.2.4) de $S$ est $S$ tout entier, donc, en particulier, que $X = \tilde{X}_\eta$ est spécial.


4. \underline{Les conclusions}

Mettant ensemble 1.1 et 3.2, on trouve :

\underline{Théorème} 4.1. \underline{Fixons une caractéristique} $p$, \underline{une dimension} $n$, \underline{et un multidegré} $(d_1,\ldots,d_r)$. \underline{Désignons par}

\begin{equation}\tag{4.1.0}
X_V \longrightarrow V
\end{equation}
\underline{l'intersection complète lisse de dimension} $n$ \underline{et multidegré} $(d_1,\ldots,d_r)$ universelle \underline{en caractéristique} $p$.

\underline{Si}

\begin{equation}\tag{4.1.1}
\sum_{1}^{r} d_i \geqslant n + r + 1 \quad ,
\end{equation}
\underline{alors il existe un ouvert dense} $U \subset V$ \underline{tel que}

\begin{equation}\tag{4.1.2}
\underline{\text{pour}}\ u \in U, \ \underline{\text{on a}}\ H^n(X_{\bar{u}}, \mathbb{Q}_\ell) \neq N^1 H^n(X_{\bar{u}}, \mathbb{Q}_\ell) \quad .
\end{equation}

%%%%%%%%%%  scan idx 390  |  folio 383  %%%%%%%%%%

Mettant ensemble 1.3 et 3.2, on trouve de meme .

\underline{Théorème} 4.2. \underline{Soit} $X$ \underline{une intersection complète,} \underline{lisse de dimension} $n \geqslant 1$ \underline{sur un corps} $k$ . \underline{Fixons un entier}

\begin{equation}\tag{4.2.1}
d \geqslant 1 + n \quad .
\end{equation}
\underline{Désignons par}

\begin{equation}\tag{4.2.2}
X.H_u \longrightarrow U
\end{equation}
\underline{la famille} universelle \underline{des sections de} $X$ \underline{par des hypersurfaces de degré} $d$, \underline{qui sont lisses de dimension} $n - 1$. \underline{Il existe un ouvert dense} $T \subset U$ \underline{tel que}

\begin{equation}\tag{4.2.3}
\underline{\text{pour}}\ t \in T, \ \underline{\text{on a}}\ H^{n-1}(X.H_{\bar{t}}, \mathbb{Q}_\ell) \neq N^1 H^{n-1}(X.H_{\bar{t}}, \mathbb{Q}_\ell) \quad .
\end{equation}

\underline{BIBLIOGRAPHIE}

[1]  J.P.  SERRE.  Faisceaux algébriques cohérents, Annals of Math, 61, 1955, pp.197-278.

[2]  J.P.  SERRE.  Sur la topologie des variétés algébriques en caractéristique $p$, dans Symposium Internacional de topologia algebraica, Mexico 1958.


%%%%%%%%%%  scan idx 391  |  folio 384  %%%%%%%%%%


§ 5. \underline{Appendice} , \underline{Théorèmes d'intégralité}. Par P. Deligne

5.0. Soient $p$ un nombre premier, $q$ une puissance de $p$, $\mathbb{F}_q$ un corps à $q$ éléments, $E$ le complété d'un corps de nombresen une place divisant un nombre premier $\ell$ distinct de $p$, et $T$ un ensemble de nombres premiers (par exemple $T = \emptyset$).

Un élément $\alpha$ d'une clôture algébrique de $E$ est dit \underline{T-entier} s'il est algébrique sur $\mathbb{Q}$ et entier sur $\mathbb{Z}\bigl[(1/r)_{r \in T}\bigr]$. Un endomorphisme $F$ d'un vectoriel de dimension finie sur $E$ est dit \underline{T-entier} si ses valeurs propres sont T-entières.

Soient $k$ un corps fini à $q'$ éléments, $\bar{k}$ une clôture algébrique de $k$ et $\mathcal{G} = \mathcal{G}\mathrm{al}(\bar{k}/k)$ le groupe de Galois de $k$. On désigne par $f \in \mathcal{G}$ la substitution de Frobenius, donnée par

\begin{equation*}
f(x) = x^{q'} \qquad ,
\end{equation*}
et par $F$ le Frobenius ``géométrique'', inverse de $f$ dans $\mathcal{G}$. Une représentation $\rho$ de $\mathcal{G}$ dans un vectoriel de dimension finie sur $E$ est dite T-entière si $\rho(F)$ est T-entier. Si $X$ est un schéma de type fini sur $\mathbb{F}_q$ et $x$ un point fermé de $X$, on désigne par $F_x$ l'élément de Frobenius géométrique du groupe de Galois de $k(x)$ ; si $\bar{x}$ est un point géométrique de $X$ localisé en $x$ et $\underline{G}$ un E-faisceau sur $X$, $F_x$ agit par transport de structure sur la fibre $G_{\bar{x}}$ .

\underline{Définition} 5.1. \underline{Un E-faisceau constructible} $\underline{G}$ \underline{sur un schéma} $X$ \underline{de type fini sur} $\mathbb{F}_q$ \underline{est dit T-entier si, pour tout point géométrique} $\bar{x}$ \underline{de} $X$ \underline{localisé en un point fermé} $x$ \underline{de} $X$, $F_x$ \underline{est un automorphisme T-entier de} $G_{\bar{x}}$ .


%%%%%%%%%%  scan idx 392  |  folio 385  %%%%%%%%%%

Si $a : X \longrightarrow \mathrm{Spec}(\mathbb{F}_q)$ est un schéma de type fini sur $\mathbb{F}_q$ et $\underline{G}$ un E-faisceau constructible sur $X$, on désignera par $\underline{H}^*(X,\underline{G})$ les E-faisceaux $R^*a_*(\underline{G})$ sur $\mathrm{Spec}(\mathbb{F}_q)_{\acute{e}t}$ ; si $\overline{\mathbb{F}}_q$ est une clôture algébrique de $\mathbb{F}_q$ et si $\overline{X}$ est le schéma déduit de $X$ par extension des scalaires de $\mathbb{F}_q$ à $\overline{\mathbb{F}}_q$, les $\underline{H}^*(X,\underline{G})$ s'identifient aux groupes de cohomologie $H^*(\overline{X},\underline{G})$, considérés comme modules galoisiens, i.e. munis d'une action du Frobenius géométrique $F$. De même en cohomologie à support propre.

\underline{Lemme} 5.2.1. \underline{Soit} $\underline{G}$ \underline{un E-faisceau constructible sur un schéma} $X$ \underline{de type fini sur} $\mathbb{F}_q$ \underline{et de dimension} $\leq n$.

(i) \underline{Il existe une partie fermée} $Y$ \underline{de dimension} $0$ \underline{de} $X$ \underline{telle que la flèche évidente de} $\underline{H}^0(X,\underline{G})$ \underline{dans} $\underline{H}^0(Y,\underline{G})$ \underline{soit injective.}

(ii) \underline{Si} $X$ \underline{est séparé, et si} $U$ \underline{est un ouvert de} $X$ \underline{dont le complémentaire} $Z$ \underline{est de dimension} $< n$, \underline{il existe une partie fermée} $Y$ \underline{de dimension} $0$ \underline{de} $U$ \underline{et une flèche surjective de} $\underline{H}^0(Y,\underline{G})(-n)$ \underline{dans} $\underline{H}^{2n}_c(X,G)$.

L'assertion (i) est évidente. Prouvons (ii). Quitte à remplacer $X$ par $X_{\mathrm{red}}$ et à rétrécir $U$, on peut supposer que $U$ est lisse et que $G|U$ est constant tordu. On a, puisque $\dim(Z) < n$,

\begin{equation*}
0 = \underline{H}^{2n-1}_c(Z,\underline{G}) \longrightarrow \underline{H}^{2n}_c(U,\underline{G}) \stackrel{\sim}{\longrightarrow} \underline{H}^{2n}_c(X,\underline{G}) \longrightarrow \underline{H}^{2n}_c(Z,G) = 0
\end{equation*}
et, par dualité de Poincaré (SGA 4 XVIII et [2])

\begin{equation*}
\underline{H}^{2n}_c(U,\underline{G}) \quad \simeq \quad \underline{H}^0(U,\underline{G}^{\vee}(n))^{\vee} \qquad .
\end{equation*}
On conclut par (i).


%%%%%%%%%%  scan idx 393  |  folio 386  %%%%%%%%%%

Théorème 5.2.2. Soient $X$ un schéma séparé de type fini sur $\mathbb{F}_q$ et de dimension $\leq n$, et $\underline{G}$ un E-faisceau constructible T-entier. Alors, les endomorphismes $F$ et $q^{n-i} F$ \underline{de} $H^i_c (X,\underline{G})$ sont T-entiers.

Le cas où $\dim X \leq 0$ est trivial, et les cas $i = 0$ et $i = 2n$ du théorème résultent de 5.2.1. Soit $Z(X,\underline{G},t)$ la série formelle

\begin{equation*}
Z(X,\underline{G},t) \;=\; \prod_{x \in X} \det (1-F_x \, t^{\deg(x)} \;;\; \underline{G}_x)^{-1} \in F[[t]] \;.
\end{equation*}
D'après (SGA 5 XV 3.2), on a

\begin{equation}\tag{5.2.3}
Z(X,\underline{G},t) \;=\; \prod_i \det (1-Ft \;;\; \underline{H}^i_c(X,\underline{G}))^{(-1)^{i+1}} \;.
\end{equation}
Prouvons 5.2.2 pour $n = 1$. On a alors $\underline{H}^i_c(X,\underline{G}) = 0$ pour $i > 2$ (SGA 4 X 4.3 ou XVII 5.2.8.1). Par hypothèse, la série formelle $Z(X,\underline{G},t)$ est à coefficients T-entiers. D'après les cas $i = 0$ et $i = 2n = 2$ de 5.2.2, le produit de $Z(X,\underline{G},t)$ par

\begin{equation*}
\det(1-Ft \;;\; \underline{H}^0_c(X,\underline{G})). \; \det(1-Ft \;;\; \underline{H}^2_c(X,\underline{G}))
\end{equation*}
est encore à coefficients T-entiers. Ce produit n'est autre que $\det(1-Ft \;;\; \underline{H}^1_c(X,\underline{G}))$, et le théorème en résulte pour $n = 1$.

Procédons maintenant par récurrence sur $n = \dim(X)$.

Soit $U$ un ouvert de Zariski de $X$, dont le complément $Z$ soit de dimension $< n$, et tel qu'il existe un morphisme $f : U \longrightarrow Y$, de but une courbe et à fibres de dimension $\leq n-1$.


%%%%%%%%%%  scan idx 394  |  folio 387  %%%%%%%%%%

\begin{center}
\begin{tikzcd}
U \arrow[r, hook, "j"] \arrow[d, "f"'] & X & Z \arrow[l, hook', "i"'] \\
Y
\end{tikzcd}
\end{center}
% STRUCTURE: 4 nodes exactly. Top row, left to right: U, X, Z. Second row, directly below U: Y. Arrows
%   (3 total): 1. U -> X : horizontal, pointing RIGHT, hooked tail at U (open immersion, ↪), label "j"
%   ABOVE the arrow. 2. Z -> X : horizontal, pointing LEFT (arrowhead at X), hooked tail at Z
%   (immersion, ↩), label "i" ABOVE the arrow. 3. U -> Y : vertical, pointing DOWN, plain arrow with
%   head at Y, label "f" to the LEFT of the arrow shaft. No other arrows; no arrow between Z and Y or X
%   and Y. Not a node: a typed full stop "." is printed on the Y baseline, well to the right of Y
%   (sentence punctuation closing the displayed diagram).
La suite exacte

\begin{equation*}
0 \longrightarrow j_! j^* \underline{G} \longrightarrow \underline{G} \longrightarrow i_! i^* \underline{G} \longrightarrow 0
\end{equation*}
induit des suites exactes

\begin{equation*}
\underline{H}^i_c(U,\underline{G}) \longrightarrow \underline{H}^i_c(X,\underline{G}) \longrightarrow \underline{H}^i_c(Z,\underline{G})
\end{equation*}
de sorte que par récurrence, il suffit de vérifier le théorème pour le faisceau $\underline{G}$ sur $U$. Grâce à l'hypothèse de récurrence appliquée aux fibres de $f$, et au théorème de changement de base pour un morphisme propre, on sait que les faisceaux $R^i f_! \underline{G}$ sur $Y$ sont T-entiers, et que les faisceaux $R^i f_! \underline{G}(i-(n-1))$ sont T-entiers. La suite spectrale de Leray pour $f$ s'écrit

\begin{equation*}
E_2^{pq} = \underline{H}^p_c(Y,R^q f_! \underline{G}) \;\Longrightarrow\; \underline{H}^{p+q}_c(U,\underline{G}) \;.
\end{equation*}
En vertu de ce qui précède, et du théorème déjà connu pour $Y$, les $E_2^{pq}$ sont T-entiers, et les $E_2^{pq}(p+q-n)$ le sont aussi. Le théorème en résulte.

\underline{Corollaire} 5.3. \underline{Soient} $X$ \underline{un schéma de type fini sur} $\mathbb{F}_q$, \underline{de dimension} $\leq 1$, $U$ \underline{un ouvert de} $X$, $j : U \hookrightarrow X$ \underline{le morphisme d'inclusion et} $\underline{G}$ \underline{un} E-\underline{faisceau constructible sur} $U$. \underline{Alors} 


%%%%%%%%%%  scan idx 395  |  folio 388  %%%%%%%%%%

(a)\quad \underline{Si} $\underline{G}$ \underline{est $T$-entier}, \underline{alors} $j_*\underline{G}$ \underline{est $T$-entier}.

(b)\quad \underline{Si pour tout} $x \in U$, $\underline{G}^\vee_x$ \underline{est $T$-entier}, \underline{alors}, \underline{pour tout} $x \in X$, $(j_*\underline{G})_x^\vee$ \underline{est $T$-entier}.

On se ramène au cas $X$ affine, puis $X$ projectif.

Définissons $\underline{H}$ par la suite exacte

\begin{equation*}
0 \longrightarrow j_!\underline{G} \longrightarrow j_*\underline{G} \longrightarrow \underline{H} \longrightarrow 0 \quad .
\end{equation*}
La suite exacte de cohomologie correspondante fournit, si $Y$ est le support (fini) de $\underline{H}$, une suite exacte

\begin{equation*}
\underline{H}^0(U,\underline{G}) \longrightarrow \underline{H}^0(Y,\underline{H}) \longrightarrow \underline{H}^1_c(U,\underline{G}) \quad .
\end{equation*}
On sait, par 5.2.1 et par le théorème 5.2 que $\underline{H}^0(U,\underline{G})$ et $\underline{H}^1_c(U,\underline{G})$ sont $T$-entiers ; il en résulte que $\underline{H}^0(Y,\underline{H})$ est $T$-entier, ce qui prouve (a).

Si $X'$ est le normalisé de la partie purement de dimension un de $X_{\mathrm{red}}$ et $U'$ un ouvert de $U$, dense dans le lieu lisse de dimension 1 de $U_{\mathrm{red}}$, et tel que $\underline{G}$ restreint à $U$ soit constant tordu :

\begin{center}
\begin{tikzcd}
U' \arrow[r, hook, "j'"'] \arrow[d, hook] & X' \arrow[d, hook, "p"] & {} \\
U \arrow[r, hook, "j"'] & X & Y \arrow[l, hook']
\end{tikzcd}
\end{center}
% STRUCTURE: Untagged square-plus-tail diagram, 5 nodes in a 2-row layout (top row 2 nodes, bottom row 3
%   nodes; the top-right cell of the third column is empty). NODES (5 total): (1) U' — top left; (2) X'
%   — top right of top row; (3) U — bottom left, directly under U'; (4) X — bottom middle, directly
%   under X'; (5) Y — bottom right. ARROWS (5 total): - A1: U' -> X'. Horizontal, pointing RIGHT.
%   Hooked/mono arrow: the ⊂-hook is drawn at the TAIL (left, at U'). Label "j'" printed BELOW the arrow
%   (in the gap between the two rows, left of centre). - A2: U' -> U. Vertical, pointing DOWN.
%   Hooked/mono arrow: the hook is drawn at the TAIL (top, at U'), i.e. a ⊂ rotated 90° (∩-shaped curl
%   opening downward, its right leg free). Plain barbed head at the bottom. No label. - A3: X' -> X.
%   Vertical, pointing DOWN. Hooked/mono arrow drawn with the SAME tail hook as A2 (∩-shaped curl at the
%   top, at X'). Plain barbed head at the bottom. Label "p" printed to the RIGHT of the arrow, at
%   mid-height. - A4: U -> X. Horizontal, pointing RIGHT. Hooked/mono arrow with the ⊂-hook at the TAIL
%   (left, at U). Label "j" printed BELOW the arrow, at centre. - A5: Y -> X. Horizontal, pointing LEFT
%   (arrowhead at X). Hooked/mono arrow with the ⊃-hook at the TAIL (right, at Y). PUNCTUATION (not a
%   node): a comma "," is printed on the baseline of the bottom row, well to the right of Y, closing the
%   displayed diagram.
alors $p_*j'_*\,\underline{G} \mid Y \simeq j_*\underline{G}|Y$. Il suffit donc d'étudier le cas où $X$ est lisse de dimension un et où $\underline{G}|U$ est constant tordu. Le faisceau $\underline{G}^\vee$ est alors $T$-entier, de sorte qu'on sait déjà que pour tout $x \in X$ $(j_* \underline{G}^\vee)_x$
%%%%%%%%%%  scan idx 396  |  folio 389  %%%%%%%%%%
est $T$-entier. Désignnons par $\mathcal{G}_x$ le groupe de Galois du corps des fractions de $\hat{\mathcal{O}}_{X,x}$ pour $x \in X$, et par $\mathfrak{I}_x$ le groupe d'inertie. Le faisceau $\underline{G}$ définit alors une représentation, encore notée $\underline{G}$, de $\mathcal{G}_x$, et le $\mathcal{G}_x/\mathfrak{I}_x = \mathcal{G}\mathrm{al}(\overline{k(x)}/k(x))$-module $(j_*\underline{G}^\vee)_x$ (resp $(j_*\underline{G})_x^\vee$) s'identifie alors au $\mathcal{G}_x/\mathfrak{I}_x$-module des invariants (resp. des coinvariants) dans $\underline{G}^\vee$. Il reste à prouver :

\underline{Lemme} 5.3.1.\quad \underline{Soient} $K$ \underline{un corps local de corps résiduel} $\mathbb{F}_q$, $G$ \underline{son groupe de Galois}, $I$ \underline{le groupe d'inertie}, \underline{et} $\rho$ \underline{une représentation continue de} $G$ \underline{dans un vectoriel} $V$ \underline{de dimension} $n$ \underline{sur} $E$ (cf. 5.0), $V^I$ \underline{l'espace des invariants sous} $I$, \underline{et} $V_I$ \underline{l'espace des coinvariants sous} $I$, \underline{considérés comme} $G/I$-\underline{modules}. \underline{Soient} $\alpha_1 \ldots \alpha_n$ \underline{les valeurs propres de} $F$ \underline{agissant sur} $V^I$. \underline{Alors, il existe des entiers} $m_i \geq 0$ \underline{tels que les valeurs propres de} $F$ \underline{agissant sur} $V_I$ \underline{soient les} $\alpha_i q^{m_i}$.

Soit $I'$ le plus grand sous-groupe de $I$ d'ordre premier à $\ell$. Le groupe $\rho(I')$ est fini, étant un sous-groupe compact d'ordre premier à $\ell$ de $\mathrm{GL}(V)$. On a donc

\begin{equation*}
V_I \quad \simeq \quad (V^{\rho(I')})_I
\end{equation*}
et, remplaçant $V$ par $V^{I'}$, on peut supposer que $I'$ agit trivialement sur $V$. On a $I/I' \simeq \mathbb{Z}_\ell(1)$ (en tant que $G/I$-modules) ; soit $T$ un générateur topologique de $I/I'$, et, pour chaque $\alpha$, soit $V^\alpha$ le sous-espace propre généralisé de $T$ dans $V$ de valeur propre $\alpha$ : on a $V = \oplus\, V^{(\alpha)}$ ; et la décomposition $V = V^{(1)} \oplus \sum_{\alpha \neq 1} V^{(\alpha)}$ est stable par $G$, et compatible
%%%%%%%%%%  scan idx 397  |  folio 390  %%%%%%%%%%
au passage au dual ; on peut donc remplacer $V$ par $V^{(1)}$, ce qui permet de supposer que $I$ agit de façon unipotente.

\underline{Lemme} 5.3.2. \underline{Soient} $V$ \underline{un vectoriel de dimension finie sur un corps} $k$ \underline{de caractéristique} $0$, $I$ \underline{un groupe additif à un paramètre}, $\rho$ $I \longrightarrow \mathrm{Gl}(V)$ \underline{une représentation de} $I$ \underline{dans} $V$, $N$ \underline{un générateur infinitésimal de} $I$ \underline{et posons}

\begin{equation*}
V^{I,k} = V^I \ \cap \ \mathrm{Im}(N^k)
\end{equation*}
\begin{equation*}
V_{I,k} = \mathrm{Im}(\mathrm{Ker}(N^{k+1}) \longrightarrow V_I) \quad .
\end{equation*}
\underline{Alors, la flèche} $N^k$ \underline{induit un isomorphisme}

\begin{equation*}
N^k : V_{I,k}/V_{I,k-1} \ \overset{\sim}{\longrightarrow} \ V^{I,k}/V^{I,k+1}
\end{equation*}
\underline{et définit des isomorphismes canoniques}

\begin{equation*}
V_{I,k}/V_{I,k-1} \otimes I^{\otimes k} \ \simeq \ V^{I,k}/V^{I,k+1}
\end{equation*}
Les représentations $V$ de $I$ s'identifient, via $N$, aux $k[T]$-modules tués par une puissance de $T$. Il suffit de vérifier 3.2 lorsque $V$ correspond à un $k[T]$-module irréductible $k[T]/T^k$, ce qui est trivial.

Achevons la démonstration de 5.3.1. Le sorite 5.3.2 fournit des filtrations croissantes et décroissantes $V_{I,k}$ et $V^{I,k}$ de $V_I$ et $V^I$, et des isomorphismes de $G/I$-modules

\begin{equation*}
V^{I,k}/V^{I,k+1} \ \simeq \ V_{I,k}/V_{I,k-1}(k) \quad ,
\end{equation*}
d'où les assertions.


%%%%%%%%%%  scan idx 398  |  folio 391  %%%%%%%%%%

\underline{Théorème} 5.4. \underline{Soient} $X$ \underline{un schéma de type fini sur} $\mathbb{F}_q$, \underline{séparé et de dimension} $n$, \underline{et} $\underline{G}$ \underline{un} $E$-\underline{faisceau constructible sur} $X$. \underline{On suppose que pour tout} $x \in X$, \underline{l'endomorphisme} $F_x$ \underline{de} $\underline{G}_x^\vee$ \underline{est} $T$-\underline{entier}. \underline{Alors l'endomorphisme} $q^i F^\vee$ \underline{de} $\underline{H}_c^i(X,G)^\vee$ \underline{est} $T$-\underline{entier}, \underline{ainsi que l'endomorphisme} $q^n F^\vee$ \underline{de} $\underline{H}_c^i(X,G)^\vee$.

On procède comme dans le théorème 5.2 pour se ramener au cas où $X$ est une courbe lisse et où $\underline{G}$ est constant tordu sur $X$. Soit $j$ l'inclusion de $X$ dans la courbe projective et lisse $X'$ qui compactifie $X$.

On sait (dualité de Poincaré) que $\underline{H}^i(X',j_*G)$ est dual de $\underline{H}^{2-i}(X',j_* G^\vee(1))$. En vertu du corollaire 5.3 (a), le théorème 5.2 s'applique à $j_*G^\vee$ :

\begin{align*}
\underline{H}^0(X',j_*G^\vee) &\ \simeq \ \underline{H}^2(X',j_*G)^\vee(-1) && \text{est \(T\)-entier,} \\
\underline{H}^1(X',j_*G^\vee) &\ \simeq \ \underline{H}^1(X',j_*G)^\vee(-1) && \text{est \(T\)-entier,} \\
\underline{H}^2(X',j_*G^\vee)(1) &\simeq \underline{H}^0(X',j_*G)^\vee && \text{est \(T\)-entier,}
\end{align*}
ce qui prouve le théorème pour $j_*G$. Définissons $H$ par la suite exacte

\begin{equation*}
0 \longrightarrow j_!G \longrightarrow j_*G \longrightarrow H \longrightarrow 0 \quad .
\end{equation*}
La suite exacte de cohomologie fournit, si $Y$ est le support de $H$,

\begin{center}
\begin{tikzcd}[row sep=large, column sep=large]
 {} & \underline{H}^0_c(X,G) \arrow[r, hook] & \underline{H}^0(X',j_*G) & , \\
\underline{H}^0(Y,H) \arrow[r] & \underline{H}^1_c(X,G) \arrow[r] & \underline{H}^1(X',j_*G) & , \\
 {} & \underline{H}^2_c(X,G) \arrow[r, "\sim"] & \underline{H}^2(X',j_*G) & ,
\end{tikzcd}
\end{center}
% STRUCTURE: Trois lignes alignées en colonnes (pas de flèches verticales ; ce n'est pas un diagramme
%   commutatif mais un fragment de suite exacte longue disposé sur trois lignes alignées). Noeuds
%   mathématiques : 7 au total. Ligne 1 (colonne 1 vide) : noeud A = \underline{H}^0_c(X,G) ; noeud B =
%   \underline{H}^0(X',j_*G). Cellule de ponctuation finale : « , ». Flèche 1 : A -> B, horizontale vers
%   la droite, à crochet (hook / mono, tracée « ⊂——→ » dans le dactylogramme), sans étiquette. Ligne 2 :
%   noeud C = \underline{H}^0(Y,H) (colonne 1, en retrait à gauche des autres lignes) ; noeud D =
%   \underline{H}^1_c(X,G) (colonne 2, aligné sous A) ; noeud E = \underline{H}^1(X',j_*G) (colonne 3,
%   aligné sous B). Cellule de ponctuation finale : « , ». Flèche 2 : C -> D, horizontale vers la
%   droite, flèche pleine simple, sans étiquette. Flèche 3 : D -> E, horizontale vers la droite, flèche
%   pleine simple, sans étiquette. Ligne 3 (colonne 1 vide) : noeud F = \underline{H}^2_c(X,G) (colonne
%   2) ; noeud G = \underline{H}^2(X',j_*G) (colonne 3). Cellule de ponctuation finale : « , ». Flèche 4
%   : F -> G, horizontale vers la droite, flèche pleine (non pointillée), étiquetée « ~ » au-dessus du
%   fût (flèche d'isomorphisme). Toutes les lettres H portent le soulignement de notation ; les
%   faisceaux G à l'intérieur des parenthèses ne sont pas soulignés.

%%%%%%%%%%  scan idx 399  |  folio 392  %%%%%%%%%%
de sorte que le théorème pour $(X,G)$ résulte du théorème pour $(X',j_*G)$.

5.5. La méthode précédente de fibration en courbes permet aussi d'obtenir des majorations à la Lang-Weil [3] sur les valeurs absolues complexes des valeurs propres de Frobenius.

Soient $X$ un schéma de type fini sur $\mathbb{F}_q$, $\underline{G}$ un E-faisceau constructible sur $X$, et $\lambda \in \mathbb{R}$. Si, pour tout point $x \in X$, avec $\sharp\, k(x) = q_x$, les valeurs propres de $F_x$ agissant sur $\underline{G}_x$ sont des nombres algébriques $\alpha$ dont toutes les valeurs absolues complexes vérifient

\begin{equation*}
|\alpha| \;\leqslant\; q_x^{\lambda} \qquad ,
\end{equation*}
on écrira

\begin{equation*}
\log |\underline{G}| \;\leqslant\; \lambda \qquad .
\end{equation*}
\underline{Lemme} 5.5.1. \underline{Soit} $X$ \underline{une courbe sur} $\mathbb{F}_q$, \underline{par quoi on entend ici un schéma} $X$ \underline{de type fini sur} $\mathbb{F}_q$, \underline{séparé et de dimension} $\leqslant 1$. \underline{Soit} $\underline{G}$ \underline{un E-faisceau constructible sur} $X$ \underline{tel que} $\log |\underline{G}| \leqslant \lambda$ \underline{avec} $\lambda \in \mathbb{R}$. \underline{Alors les valeurs absolues complexes des valeurs propres} $\alpha$ \underline{de} $F$ \underline{agissant sur} $\underline{H}^1_c(X,G)$ \underline{vérifient} $|\alpha| \leqslant q^{\lambda+1}$.

Soit $\sigma$ un plongement complexe de E, ou simplement un plongement complexe de la fermeture algébrique de $\mathbb{Q}$ dans E.

Il suffit de vérifier que, quel que soit $\sigma$, les zéros $\beta$ de la fraction rationnelle $\sigma Z(X,\underline{G},t)$ vérifient $|\beta^{-1}| \leqslant q^{\lambda+1}$. On sait déjà en
%%%%%%%%%%  scan idx 400  |  folio 393  %%%%%%%%%%
effet par 5.2.1 si $\alpha$ est une valeur propre de Frobenius agissant sur $\underline{H}^0_c(X,\underline{G})$ ou sur $\underline{H}^2_c(X,\underline{G})$, on a $|\sigma \alpha| \leqslant q^{\lambda+1}$.

Vérifions que la série

\begin{equation*}
\log \sigma Z(X,\underline{G},t) \;=\; \sum_k \frac{t^k}{k} . \sum_{x \in X(\mathbb{F}_{q^k})} \sigma \; \mathrm{Tr}(F,x^*\underline{G})
\end{equation*}
est convergente pour $|t| < q^{-\lambda-1}$ . Représentant $X$ comme revêtement ramifié de $\mathbb{P}^1$, on obtient trivialement une majoration

\begin{equation*}
\sharp \; X(\mathbb{F}_{q^k}) \;\leqslant\; A.q^k \qquad .
\end{equation*}
Par hypothèse, on a d'autre part

\begin{equation*}
\left| \; \sigma \; \mathrm{Tr}(F,x^*G) \right| \;\leqslant\; B.(q^k)^{\lambda} \qquad .
\end{equation*}
Le coefficient $c_k$ de $t^k$ dans $\sigma Z (X,\underline{G},t)$ vérifie donc

\begin{equation*}
\left| \, c_k \right| \;\leqslant\; \frac{1}{k} \;.\; C. \; (q^{1+\lambda})^k
\end{equation*}
et l'assertion en résulte.

\underline{Théorème} 5.5.2. \underline{Soient} $X$ \underline{un schéma de type fini sur} $\mathbb{F}_q$, \underline{séparé et de dimension} $\leqslant n$, \underline{et} $G$ \underline{un E-faisceau constructible sur} $X$ \underline{tel que} $\log|G| \leqslant \lambda$ \underline{avec} $\lambda \in \mathbb{R}$ .

\underline{Alors, les valeurs absolues complexes des valeurs propres} $\alpha$ \underline{de Frobenius agissant sur} $H^i_c(X,\underline{G})$ \underline{vérifient}

\begin{equation*}
|\alpha| \;\leqslant\; q^{\lambda+i} \qquad \underline{\mbox{et}}
\end{equation*}
\begin{equation*}
|\alpha| \;\leqslant\; q^{\lambda+n} \qquad .
\end{equation*}

%%%%%%%%%%  scan idx 401  |  folio 394  %%%%%%%%%%

Si $n \leq 0$, 5.5.2 est trivial. Si $n = 1$, les cas $i = 0$ et $i = 2$ de 5.5.2 résultent de 5.2.1, tandis que le cas $i = 1$ résulte de 5.5.1. On procède alors par récurrence sur $n$, comme dans la démonstration de 5.5.2.

\underline{Corollaire} 5.5.3. \underline{Soit} $X$ \underline{un schéma séparé de type fini de dimension} $n$ \underline{sur} $\mathbb{F}_q$ . \underline{Alors, les valeurs propres} $\alpha$ \underline{de l'endomorphisme de Frobenius géométrique} $F$ \underline{de} $\underline{H}^i_c(X,\mathbb{Q}_\ell)$ \underline{sont des entiers algébriques}. \underline{De plus}

(i)\quad \underline{Pour} $\ell' \neq p$, $\alpha$ \underline{est une unité} $\ell'$\underline{-adique}.

(ii)\quad \underline{Pour} $i > n$, $\alpha$ \underline{est divisible par} $q^{i-n}$.

(iii)\quad \underline{Pour} $i \leq n$, (resp. \underline{pour} $i \geq n$), $q^i \alpha^{-1}$ (resp. $q^n \, \alpha^{-1}$) \underline{est un entier algébrique}.

(iv)\quad \underline{Pour} $0 < i$, \underline{les valeurs absolues complexes de} $\alpha$ \underline{vérifient}

\begin{equation*}
|\alpha| \leq q^{i-\frac{1}{2}} \qquad \mathrm{et} \qquad |\alpha| \leq q^n \quad .
\end{equation*}
D'après 5.2.2, $\alpha$ et $q^{n-i} \, \alpha$ sont des entiers algébriques. Les assertions (i) à (iii) résultent alors de 5.4. Pour $i > n$, (iv) résulte de 5.5.2. Pour obtenir (iv) quand $i \leq n$, il va falloir conjuguer 5.5.2 à l'hypothèse de Riemann pour les courbes. Procédons par récurrence sur $n$, le cas $n = 0$ étant trivial. Si $U$ est un ouvert de $X$, de complémentaire $Z$ de dimension $< n$, la suite exacte

\begin{equation*}
\underline{H}^i_c(U,\mathbb{Q}_\ell) \longrightarrow \underline{H}^i_c(X,\mathbb{Q}_\ell) \longrightarrow \underline{H}^i_c(Z,\mathbb{Q}_\ell)
\end{equation*}
montre qu'il suffit de prouver 5.2.2 pour $U$. Quitte à remplacer au préalable $X$ par $X_{\mathrm{red}}$, ceci nous ramène au cas où $X$ est lisse affine,
%%%%%%%%%%  scan idx 402  |  folio 395  %%%%%%%%%%
et tel qu'il existe un morphisme lisse $f : X \longrightarrow S$ vérifiant les conditions suivantes :

(a)\quad $\dim(S) \leq n-1$

(b)\quad les fibres de $f$ sont purement de dimension 1.

Puisque les fibres de $f$ sont des courbes affines, on a

\begin{equation*}
R^0 f_! \; \mathbb{Q}_\ell = 0 \quad .
\end{equation*}
D'après Weil [4], on a

\begin{equation}\tag{5.5.3.1}
\log \, \bigl| R^1 f_! \; \mathbb{Q}_\ell \bigr| \leq 1/2 \quad ,
\end{equation}
et d'après 5.2.1 appliqué aux fibres de $f$, on a

\begin{equation}\tag{5.5.3.2}
\log \, \bigl| R^2 f_! \; \mathbb{Q}_\ell \bigr| \leq 1 \quad .
\end{equation}
La suite spectrale de Leray pour $f$ se réduit ici à une suite exacte longue

\begin{equation*}
\underline{H}^{i-2}_c(S,R^2 f_! \mathbb{Q}_\ell) \longrightarrow \underline{H}^i_c(X,\mathbb{Q}_\ell) \longrightarrow \underline{H}^{i-1}_c(S,R^1 f_! \mathbb{Q}_\ell) \qquad .
\end{equation*}
D'après (5.5.3.1) (resp. 5.5.3.2) et 5.5.2, les valeurs absolues complexes des valeurs propres $\alpha$ de Frobenius agissant sur $\underline{H}^{i-1}_c(S,R^1 f_! \mathbb{Q}_\ell)$ (resp. sur $\underline{H}^{i-2}_c(S,R^2 f_! \mathbb{Q}_\ell)$) vérifient

\begin{equation*}
|\alpha| \leq q^{i-1+\frac{1}{2}} \quad = q^{i-\frac{1}{2}}
\end{equation*}
\begin{equation*}
(\mathrm{resp.} \; |\alpha| \leq q^{i-2+1} \quad = q^{i-1} \leq q^{i-\frac{1}{2}} \qquad .
\end{equation*}
Le corollaire en résulte.


%%%%%%%%%%  scan idx 403  |  folio 396  %%%%%%%%%%

Pour obtenir quelques résultats analogues aux précédents, en cohomologie sans supports, il semble nécessaire de disposer de la résolution des singularités. En dimension $\leq 1$, la solution du problème est fournie par le corollaire 5.3.

\underline{Théorème} 5.6. \underline{Soit} $N$ \underline{un entier, et supposons que l'on dispose de la résolution des singularités, sous la forme forte de Hironaka, pour les schémas de dimension} $\leq N$ \underline{de type fini sur} $\mathbb{F}_q$ \underline{(jusqu'ici, on peut donc prendre} $N = 2$ [1]). \underline{Soient} $f : X \longrightarrow Y$ \underline{un morphisme de schémas de type fini sur} $\mathbb{F}_q$ \underline{et} $\underline{G}$ \underline{un E-faisceau constructible T-entier sur} $X$. \underline{Alors, si} $\dim(X) \leq N$, \underline{les E-faisceaux} $R^i f_*\underline{G}$ \underline{sur} $Y$ \underline{sont T-entiers.}

La démonstration procède par récurrence sur la dimension $n \leq N$ de $X$. Soit $j : U \longrightarrow X$ un ouvert de $X$, lisse, affine, tel que $\underline{G}$ restreint à $U$ soit constant tordu et tel que la dimension de son complément $Z$ soit $< n$. Utilisant l'hypothèse de résolution, on sait que les $R^i j_*(j^*\underline{G})$ sont constructibles (SGA 5 I \quad , [2]); pour $i > 0$, ils sont concentrés sur $Z$ et, pour $i = 0$, $j_*(j^*G)$ ne différe de $\underline{G}$ que sur $Z$. Il en résulte, d'après l'hypothèse de récurrence, qu'il suffit de vérifier le théorème pour le faisceau $j^*G$ et les morphisme $j$ et $f\,j$. Il suffit de vérifier le théorème sous l'hypothèse additionnelle que $X$ soit affine lisse et $\underline{G}$ constant tordu (i.e. fourni par une représentation continue du groupe fondamental de $X$ dans un E-vectoriel de dimension finie).


%%%%%%%%%%  scan idx 404  |  folio 397  %%%%%%%%%%

Soient $R$ l'anneau des entiers de $E$, $\lambda$ une uniformisante, $e$ l'indice de ramification absolu et $b$ un entier $> \frac{e}{\ell-1}$ . Puisque $X$ est lisse, donc normal et à groupe fondamental profini, il existe un R-faisceau $\underline{G}^o$ constant tordu tel que $\underline{G} \simeq \underline{G}^o \otimes_R E$. Soit $p : X' \longrightarrow X$ le revêtement étale de $X$ qui trivialise le faisceau localement constant $\underline{G}^o \otimes_R R/\lambda^b$. Le faisceau $\underline{G}$ est facteur direct de $p_*p^*\underline{G}$, de sorte qu'il suffit de vérifier le théorème pour $p^*\underline{G}$ et $f \circ p$. Il reste donc à traiter le cas où les hypothèses suivantes sont vérifiées. :

(*)  $X$ est affine et lisse, et il existe un R-faisceau $\underline{G}^o$, constant tordu, tel que $\underline{G} \simeq \underline{G}^o \otimes_R E$ et que le faisceau $\underline{G}^o \otimes_R R/\lambda^b$ soit constant .

Factorisons le morphisme $f$ en un morphisme propre et une immersion ouverte $f = \bar{f} j$ :

\begin{center}
\begin{tikzcd}
X \arrow[r, hook, "j"] \arrow[d, "f"'] & X' \arrow[dl, "\bar{f}"] \\
Y &
{} \end{tikzcd}
\end{center}
% STRUCTURE: 3 nodes, 3 arrows. Nodes: X (top left), X' (top right), Y (bottom left); the bottom-right
%   position is empty. Arrows: (1) X --> X', horizontal, pointing right, hook/open-immersion arrow (the
%   tail at X is a left hook "⊂"), single arrowhead at X', label "j" placed below the middle of the
%   arrow; (2) X --> Y, vertical, pointing down, plain solid arrow with single arrowhead at Y, label "f"
%   placed to the left of the middle of the arrow; (3) X' --> Y, diagonal, running from upper right to
%   lower left, plain solid arrow with single arrowhead at Y, label "$\bar{f}$" (f with an overbar)
%   placed to the right of the middle of the diagonal. No dashed, no double, no iso-tilde, no two-headed
%   arrows; the square is not filled and carries no commutativity mark. A stray full stop is typed to
%   the right of Y, on Y's line, terminating the displayed sentence; it is punctuation, not a node, and
%   is not included in the tikzcd.
Le théorème a déjà été établi pour $\bar{f}$ (5.2), de sorte qu'il suffit de l'établir pour $j$. D'après la résolution des singularités, appliquée à $X'$, on pouvait supposer que $X'$ soit lisse et que $X$ soit le complément d'un diviseur à croisements normaux de $X'$, réunion de diviseurs lisses. L'hypothèse (*) implique que le faisceau $\underline{G}$ est modérément ramifié le long de ce diviseur. On a alors :


%%%%%%%%%%  scan idx 405  |  folio 398  %%%%%%%%%%

\underline{Lemme} 5.6.1. \underline{Soient} $X$ \underline{un schéma lisse de type fini sur} $\mathbb{F}_q$, $D$ \underline{un diviseur à croisements normaux dans} $X$, \underline{réunion de diviseurs lisses} $D_i$, $j : U \hookrightarrow X$ \underline{le complément de} $D$ \underline{dans} $X$ \underline{et} $\underline{G}$ \underline{un} $E$\underline{-faisceau constant tordu sur} $U$, \underline{modérément ramifié le long de} $D$. \underline{Alors, si} $\underline{G}$ \underline{est} $T$\underline{-entier, les faisceaux} $R^i j_*\underline{G}$ \underline{sont} $T$\underline{-entiers.}

On admettra (cf. SGA 1 XIII\hspace{2.5em}) :

\underline{Lemme} 5.6.2. (i) \underline{Sous les hypothèses} 5.6.1, \underline{soient} $D'$ \underline{une composante de} $D$, $D''$ \underline{la réunion des autres composantes,} $U' = X-D''$ \underline{et} $D^{\prime o} = D'\cap U'$ :

\begin{center}
\begin{tikzcd}
{} & D^{\prime o} \arrow[r, hook, "j'_D", swap] \arrow[d, "i^{o}"] & D' \arrow[d, "i"] \\
U \arrow[r, hook, "j_1", swap] & U' \arrow[r, hook, "j'", swap] & X
\end{tikzcd}
\end{center}
% STRUCTURE: 5 nodes on a 2-row x 3-column grid. Row 1: column 1 EMPTY; column 2 = D'^o (D-prime with
%   raised circle superscript); column 3 = D'. Row 2: column 1 = U; column 2 = U'; column 3 = X. 5
%   arrows: (1) D'^o -> D' horizontal rightward, HOOK arrow (open ⊂ tail then long arrow, i.e.
%   mono/inclusion), label j'_D printed BELOW the arrow; (2) D'^o -> U' vertical downward, plain long
%   arrow with open head, label i^o printed to the RIGHT of the arrow; (3) D' -> X vertical downward,
%   plain long arrow with open head, label i printed to the RIGHT of the arrow; (4) U -> U' horizontal
%   rightward, HOOK arrow (⊂ tail), label j_1 printed BELOW the arrow; (5) U' -> X horizontal rightward,
%   HOOK arrow (⊂ tail), label j' printed BELOW the arrow. No dashed, no double-headed, no iso-labelled
%   arrows. A full stop '.' is printed on the baseline of the bottom row, well to the right of X
%   (transcribed as the separate prose block that follows).
.

\underline{Alors, la flèche de changement de base}

\begin{equation*}
i^* Rj'_* (Rj_{1*} \underline{G}) \longrightarrow Rj'_{D*}\, i^{o*} (Rj_{1*}\underline{G})
\end{equation*}
\underline{est un isomorphisme, et les faisceaux} $i^{o*}Rj_{1*}\underline{G}$ \underline{sont constant tordus, modérément ramifiés le long de} $D' - D^{\prime o}$.

(ii) \underline{Si} $C$ \underline{est une courbe lisse sur} $X$, \underline{transversale à} $D$ \underline{en un point lisse} $x$ \underline{de} $D$,

\begin{center}
\begin{tikzcd}
C \cap U \arrow[r, "j_c", swap] \arrow[d, "k'"] & C \arrow[d, "k"] \\
U \arrow[r, "j", swap] & X
\end{tikzcd}
\end{center}
% STRUCTURE: 4 nodes on a 2-row x 2-column grid. Row 1: column 1 = C ∩ U (the letters C, cap-symbol, U
%   set with spaces on the typewriter); column 2 = C. Row 2: column 1 = U; column 2 = X. 4 arrows: (1) C
%   ∩ U -> C horizontal rightward, plain long arrow with open head (NO hook), label j_c printed BELOW
%   the arrow (subscript c is lowercase); (2) C ∩ U -> U vertical downward, plain long arrow with open
%   head, label k' printed to the RIGHT of the arrow; (3) C -> X vertical downward, plain long arrow
%   with open head, label k printed to the RIGHT of the arrow; (4) U -> X horizontal rightward, plain
%   long arrow with open head (NO hook), label j printed BELOW the arrow. No dashed, no hook, no
%   double-headed, no iso-labelled arrows. A comma ',' is printed on the baseline of the bottom row,
%   well to the right of X (transcribed as the separate prose block that follows).
,
%%%%%%%%%%  scan idx 406  |  folio 399  %%%%%%%%%%
\underline{alors, la flèche de changement de base}

\begin{equation*}
k^* Rj_*\underline{G} \longrightarrow Rj_{c*}\, k^{\prime *}\underline{G}
\end{equation*}
\underline{est un isomorphisme en} $x$.

Déduisons 5.6.1 de 5.6.2 en raisonnant par récurrence sur la dimension, et en admettant provisoirement 5.6.1 dans le cas des courbes. Il résulte de 5.6.2 (ii) et de cette hypothèse, que les faisceaux $R^i j_*\underline{G}$ sont $T$-entiers en tout point lisse de $D$. En particulier, avec les notations de 5.6.2 (i), les faisceaux $R^i j_{1*}\underline{G}$ sont $T$-entiers, de sorte que, d'après 6.2 (i) et l'hypothèse de récurrence appliquée à $i^{o*} Rj_{1*}\underline{G}$ et à $j'_D$ , les faisceaux $i^* R^i j_*\underline{G}$ sont $T$-entiers, quelle que soit la composante choisie $D'$ de $D$. Ceci prouve la conclusion de 5.6.1.

Reste le cas des courbes. Pour une courbe, sous les hypothèses 5.6.1, si $x$ est un point de $X - U$, si $\mathcal{G}$ est le groupe de Galois du corps des fractions de $\hat{\mathcal{O}}_{x,X}$ et $\mathfrak{J}$ le groupe d'inertie, si on désigne encore par $\underline{G}$ la représentation de $\mathcal{G}$ définie par $\underline{G}$ et si $\bar{x}$ est un point géométrique de $X$ localisé en $x$, on a des isomorphismes de $\mathcal{G}/\mathfrak{J}$-modules galoisiens :

\begin{equation*}
(j_*\underline{G})_{\bar{x}} \;\simeq\; \underline{G}^{\mathfrak{J}}
\end{equation*}
\begin{equation*}
(R^1 j_*\underline{G})_{\bar{x}} \;\simeq\; \underline{G}_{\mathfrak{J}}(-1)
\end{equation*}
\begin{equation*}
(R^i j_*\underline{G})_{\bar{x}} = 0 \ \text{si} \ i \geq 2
\end{equation*}
et l'assertion résulte de 5.3 (a) et 5.3.1.
